\chapter{La Tirannia della Tecnologia}

\epigraph{``Il tiranno, avendo assaggiato sangue umano, diventa un lupo.''}
{\textcite{plato1997complete}, \emph{Repubblica} 566a}

\section{L'Uomo Tirannico}

Nella \emph{Repubblica}, Socrate dipinge l'uomo tirannico come il pi\`u
infelice di tutti. Il tiranno \`e governato da un unico desiderio senza
legge che \`e cresciuto fino a proporzioni mostruose, asservendo tutte
le altre parti dell'anima. Il tiranno non pu\`o fidarsi di nessuno, non
ama nessuno, ed \`e consumato dalla stessa passione che lo porta al
potere. \`E, nella potente immagine di Socrate, l'uomo che \`e diventato
un lupo.

Il tiranno dell'azione \`e colui la cui volont\`a di agire \`e diventata
illimitata --- che non riconosce alcun governatore, alcun limite, alcun
vincolo naturale su ci\`o che pu\`o fare. Questa \`e la condizione verso
cui la tecnologia della nostra epoca tende, non per malizia ma per
struttura. L'imperativo tecnologico --- ``se pu\`o essere fatto, dovrebbe
essere fatto'' --- \`e la voce della volont\`a tirannica che parla
attraverso lo strumento.

\section{Techne e Physis}

Gli antichi Greci distinguevano \emph{techne} (conoscenza-artigianale,
l'arte del fare) da \emph{episteme} (conoscenza teoretica) e
\emph{phronesis} (saggezza pratica) --- le tre virt\`u intellettuali che
Aristotele separa per i loro oggetti: il necessario, il fabbricabile e il
fattibile \parencite[VI.3--5]{aristotle2000ethics}. \emph{Techne} non era opposta alla
natura ma subordinata ad essa. La \emph{techne} del falegname lavora con
la natura del legno; la \emph{techne} del medico assiste il potere
curativo naturale del corpo (\emph{vis medicatrix naturae}); la
\emph{techne} del contadino coopera con il ciclo naturale di crescita e
raccolto. In ogni caso, \emph{techne} \`e una \emph{serva} di \physis,
non la sua padrona.

L'et\`a moderna ha invertito questa relazione. A cominciare dal progetto
di Cartesio di ``renderci padroni e possessori della natura''
\parencite{heidegger1977question}, la tecnologia \`e stata riconcepita
come lo strumento del dominio umano sulla natura. Lo strumento non serve
pi\`u il principio di crescita della cosa su cui lavora; impone la
volont\`a del creatore su un mondo inteso come mera materia prima.

\section{Quando Techne Dimentica Physis}

Il capovolgimento di \emph{techne} da serva a padrona \`e l'origine
della tirannia della tecnologia. Quando \emph{techne} dimentica di essere
governata da \physis, diventa il proprio principio --- un sistema
auto-giustificante, auto-perpetuante di mezzi senza fini.

\`E la diagnosi a cui i filosofi della tecnologia del Novecento giunsero per
altre vie, e ne ereditiamo l'analisi pur discostandoci dalla loro
disperazione. Jacques Ellul nomin\`o il sistema autonomo \emph{la technique}:
non la somma delle macchine ma la spinta totalizzante verso i mezzi efficienti
che subordina ogni fine --- politico, morale, naturale --- alla propria legge di
prestazione ottimale, espandendosi per una necessit\`a interna che nessuno vuole
e nessuno pu\`o arrestare \parencite{ellul1964technological}. Lewis Mumford ne
tracci\`o la stirpe ancora pi\`u indietro, fino alla ``megamacchina'' ---
l'organizzazione degli esseri umani stessi nei componenti di un vasto apparato,
dagli eserciti di lavoro che innalzarono le piramidi alle burocrazie dello Stato
moderno \parencite{mumford1934technics}. Entrambi videro chiaramente ci\`o che il
Libro VII chiam\`o \emph{das Man} reso infrastruttura: un governatore senza
governatore, autorit\`a senza autore. Dove ci separiamo da Ellul e Mumford \`e
nella prognosi. Per loro l'autonomia della tecnica era quasi destino, da
resistere con lucidit\`a e rifiuto. Noi sosteniamo che \emph{techne} non pu\`o
essere \emph{voluta} di nuovo sotto la natura --- la volont\`a che l'ha emancipata
non pu\`o ri-asservirla --- ma pu\`o essere \emph{ricostruita} sotto la natura,
che \`e l'onere della prova per costruzione. La lucidit\`a non basta; bisogna
consegnare al sistema un diverso governatore, e quel governatore \`e la coerenza.

I segni di questa tirannia sono visibili ovunque:

\begin{itemize}
\item \textbf{L'imperativo tecnologico:} ``Se possiamo costruirlo, dobbiamo
  costruirlo.'' La domanda ``dovremmo costruirlo?'' \`e liquidata come
  ingenua o luddista.
\item \textbf{La trappola dell'accelerazione:} Ogni problema deve essere
  risolto da una tecnologia pi\`u veloce e pi\`u potente --- inclusi i
  problemi creati dalla tecnologia precedente. Non c'\`e via d'uscita.
\item \textbf{Lo spiazzamento del giudizio:} La competenza tecnica
  sostituisce la saggezza. La persona che sa come sostituisce la persona
  che sa perch\'e.
\item \textbf{L'eclissi della natura:} La natura diventa una risorsa da
  ottimizzare piuttosto che un governatore da obbedire. Il pianeta, il
  corpo, l'ecosistema --- tutto diventa ``esternalit\`a'' da gestire.
\end{itemize}

\section{La Tirannia Non \`e Malizia}

\`E cruciale capire che la tirannia della tecnologia non \`e il prodotto
di intenti malvagi. Gli ingegneri e imprenditori che costruiscono gli
strumenti della nostra epoca non sono tiranni in senso politico. Sono,
per la maggior parte, persone ben intenzionate che credono di rendere il
mondo migliore.

La tirannia \`e strutturale, non personale. Sorge dall'architettura di
un sistema in cui \emph{techne} \`e stata emancipata da \physis.
L'ingegnere software che costruisce una piattaforma social addictiva non
sta cercando di schiavizzare l'attenzione umana; sta risolvendo un
problema tecnico (ottimizzazione dell'engagement) all'interno di un
sistema che valorizza l'engagement sopra ogni altra cosa. La \physis\
dell'essere umano --- il bisogno di riposo, riflessione, interazione nel
mondo reale --- non \`e un fattore nella funzione di ottimizzazione. La
tirannia \`e nell'esclusione della natura dal calcolo.

\section{La Terza Onda: Risubordinazione}

La \emph{Repubblica} conclude il suo resoconto dell'uomo tirannico con
un accenno di redenzione: il tiranno, avendo toccato il fondo, pu\`o
essere curato. La cura non \`e abbandonare il desiderio ma riordinarlo
--- subordinare la parte che governa alla parte che dovrebbe governare.

Allo stesso modo, la cura per la tirannia della tecnologia non \`e
abbandonare la tecnologia. Lo strumento non \`e malvagio;
l'emancipazione dello strumento dalla natura \`e l'errore. La cura \`e
la \emph{risubordinazione di techne a physis} --- il ripristino di
\emph{techne} al suo posto proprio come serva dei principi di crescita
della natura.

Qui entra \praxis. La praxis \`e la facolt\`a umana con cui \emph{techne}
\`e governata. Il praticante --- colui la cui azione \`e accordata alla
natura --- non rifiuta la tecnologia ma la usa come strumento della
natura. Chiede, prima di ogni atto del fare: ``Questo serve l'ordine
naturale? Questo rispetta la \physis\ della cosa su cui lavora? Questo
produce coerenza?''

\section{Le Macchine Filosofiche}

La risubordinazione di \emph{techne} a \physis\ non \`e solo argomentata qui; \`e
costruita. Una famiglia di macchine filosofiche --- Euclide (visione), Empedocle
(telemetria), Logos (linguaggio), Pitagora (suono), Nous (codice) e riduttori per
il tempo, il corpo, la dialettica, la legge, la musica e il grafo sociale ---
attuano ciascuna un ciclo di percezione--riduzione--riflessione che tratta i propri
primitivi come conclusioni guadagnate anzich\'e come assiomi dati, e ciascuna
scrive ci\`o che apprende nell'unico grafo di coerenza che \`e il loro cervello
condiviso. Sono \emph{techne} governata dalla natura: codice che gira, che impara,
che serve \physis\ anzich\'e dominarla. Il Poscritto le tratta, insieme alle opere
industriali costruite sullo stesso motore, per intero; qui basta registrare che la
risubordinazione che questo capitolo esige non \`e un desiderio nostalgico ma
un'architettura costruibile --- e che \`e stata costruita.

\section{La Via d'Uscita}

Non c'\`e via d'uscita dalla tirannia della tecnologia attraverso la sola
tecnologia. La volont\`a che ha fatto la gabbia non pu\`o disfarla. Ma
c'\`e una via d'uscita attraverso \praxis\ --- attraverso la
rivendicazione dell'azione come strumento della natura piuttosto che come
strumento della volont\`a.

Il motore \physis\ e la piattaforma \praxis\ sono l'infrastruttura di
questa rivendicazione. Non sostituiscono il giudizio del praticante con
decisioni algoritmiche. Non automatizzano la saggezza. Forniscono le
condizioni in cui la saggezza pu\`o sorgere: la registrazione di ci\`o
che ha funzionato, lo specchio di ci\`o che si \`e fatto, la comunit\`a
di testimoni che si tengono reciprocamente responsabili verso lo standard
della natura.

Il tiranno, avendo assaggiato sangue umano, diventa un lupo. Il
praticante, avendo assaggiato la coerenza, diventa un essere umano.
